\begin{align*}
    &\text{Defino, para la luz: } \breve{\lambda}=\frac{\lambda M c}{a}\; \text{(parámetro)}\; \; \;
    \breve{E}=\frac{E}{Mc^2}\; \; \; \breve{L}=\frac{L}{Mca};\; \; \; c=1\; \; \; a=1\\
    &\text{A partir de ahora, no usaré el símbolo $\smallsmile$  para denotar las unidades adimensionales.}\\
    &\frac{dt}{d\lambda}= \frac{E}{\left(1-\frac{1}{r}\right)}\; \; \; \; \frac{d\phi}{d\lambda}=\frac{L}{r^2}\;\Rightarrow \;\varepsilon=\left(\frac{dr}{d\lambda} \right)^2 + \frac{L^2}{r^2}\left(1-\frac{1}{r}\right)\; \; \; \, 
    \varepsilon\equiv E^2\\
    &\text{Potencial efectivo: }V_{\text{eff}}=\frac{L^2}{r^2}-\frac{L^2}{r^3}\; \; \; \; 
\left|
\begin{array}{l}
 r=3/2 \;  \; (\text{máximo}) \; \;\forall\;  L\neq 0
 \end{array}
\right.
\end{align*}
